% ===== 原文 PDF pp.68--71（印刷 pp.67--70）=====
% tr-06 / chunk b：4.5/例 4.7 的收尾（p.68 上部）、§4.7 AGT 对应、§4.8 与 §4.8.1、表 4（p.70）、
%                    §4.8.2 非正则穿刺点起（p.71）
% 本文件 tag：4.26 与 4.27（p.68）、4.28（p.69）；表 4（p.70）

% [边界] 例 4.7 的最后一句（「…这第二种构造使 Witten 的反常 … 变得显然 [56,213]。」）
% 起于 p.67 末、收尾于 p.68 首行「anomaly [212] …」，已由 tr-06 在 tr06a.tex 中写完；
% 本文件自 p.68 的下一个自然段起，不重译。

% [原文 p.67]
一般而言，多种 Class S 构造可以对同一个 SCFT 给出互补的描述；特别地，使用非正则
穿刺点的构造对于把对称性增强以及对称性的细节（例如它们的 't Hooft 反常）显现出来
很有用。我们将在 §4.8.2 中正式引入非正则穿刺点，并在 §4.8.3 中引入它们的扭曲版本。

\SUB{4.7}{AGT 对应关系一瞥}

Alday--Gaiotto--Tachikawa（AGT）对应是 [214] 中提出的一个猜想，它把由 $\fg$ 型
6d $\cN=(2,0)$ 理论紧化得到的 4d $\cT(\fg,\cC,p)$ Class S 理论的 Nekrasov 配分函数
[76,77]（另见 [215,216]）与 $\fg$-Toda 理论的共形块 [217] 联系起来。下面我们集中
讨论 $A_1$ 型的 Class S 构造，此时 Toda 共形场论就是众所周知的 Liouville 理论
（关于这一主题的近期综述见 [218--220]）。

更确切地说，$\cT(A_1,\cC,p)$ 理论被放在一个尺寸为 $R$、压扁参数（squashing
parameter）为 $b$ 的压扁球面 $S^4_b$ 上，而 $b$ 由
\begin{equation}
  b^2(x_1^2+x_2^2) + \frac{1}{b^2}(x_3^2+x_4^2) + x_5^2 = R^2\,,
  \tag{4.26}
\end{equation}
定义，其中 $x_i\in\R$。只要该 4d 理论有 $\cN=2$ 的拉格朗日量描述，$S^4_b$ 上的
超对称配分函数就可以用局域化方法 [78,221] 显式算出（详尽综述另见 [79,222]）。
超对称局域化把路径积分约化为一个有限维矩阵模型，后者给出对矢量多重态标量实部的
一个普通积分，其被积函数由经典作用量的指数、理论中各场给出的单圈行列式加权，
并进一步被 $S^4_b$ 两个极点处点状瞬子的贡献所修饰。完整的配分函数形如
\begin{equation}
  Z_{S^4_b}(q,\bar q) = \int \dd a\,
  Z_{\mathrm{cl}}(a,q,\bar q)\, Z_{1\text{-loop}}(a)\,
  Z_{\mathrm{inst}}(a,q)\, Z_{\mathrm{inst}}(a,\bar q)\,.
  \tag{4.27}
\end{equation}
贡献 $Z_{\mathrm{cl}}$ 用经典作用量给局域化组态加权，而单圈涨落由
$Z_{1\text{-loop}}$ 计入。最后，$Z_{\mathrm{inst}}$ 表示 Nekrasov 瞬子配分函数。
\footnote{一般的秩一 4d $\cN=2$ SCFT 不是拉格朗日理论，所以可能没有已知的球面配分
函数局域化公式。这种情况发生在由非正则穿刺点构造的 Argyres--Douglas 理论中。
尽管如此，我们仍然预期其配分函数具有 (4.27) 那样、对其库仑分支某个实切片的积分
形式，其中各个构件在 Liouville（以及 Toda）共形场论中有如表 4 所示的自然解释
（例如参见 [130]）。}

% [Tutor Note: (4.26) 是「压扁」四维球面：b = 1 时回到标准 S^4（半径 R）。
% 之所以要压扁，是因为 Nekrasov 配分函数的两个 Ω 背景参数 ε_1, ε_2 正好对应
% b 与 1/b（见下面表 4 第一行），压扁球面上的局域化才能同时保留两个参数。]
% [Tutor Note: (4.27) 的四项分别对应：经典作用量、单圈行列式、以及在 S^4_b 的
% 南北两个极点处各插入一份瞬子（分别按 q 与 \bar q 展开）。这就是 Nekrasov
% 配分函数 Z_Nekrasov = Z_inst(a,q) Z_inst(a,\bar q) 的来由；在 AGT 中，
% 它对应 Liouville 共形块的「两个展开方向」。]

% [原文 p.68]
这样的配分函数被猜想等于 Liouville 共形场论（更高秩情形下是其推广 Toda 共形场论）
中某些顶点算子的关联函数。这些顶点算子插入在每个穿刺点 $z_i$ 处，并且依赖于
Class S 构造中穿刺点的类型，
\begin{equation}
  Z_{S^4_b}\big(\cT(A_1,\cC,p)\big) =
  \big\langle V_{p_1}(z_1)\cdots V_{p_n}(z_n)\big\rangle
  ^{\text{Liouville}}_{\bar{\cC}}\,,
  \tag{4.28}
\end{equation}
其中 $\bar{\cC}$ 是去掉点 $z_1,\dots,z_n$ 之后的紧黎曼面。我们不打算深入这一对应的
细节，而是把 Liouville 共形场论中的可观测量与来自 Class S 构造的秩一 SCFT 之间的
字典列在表 4 中（取自 [214]）。Liouville 共形场论中相关要素的简要回顾可参见 [218]。
关于 AGT 对应的细节与推广，我们还可以参考详尽的综述 [64] 及其详细的文献列表。

% [Tutor Note: (4.28) 里的 \bar C 与 C 的区别：C 是带穿刺点的 UV 曲线（作为 Gaiotto
% 曲线），而 \bar C 是把穿刺点真正挖掉之后剩下的紧黎曼面——Liouville 关联函数就
% 定义在它上面，z_i 是那些被挖掉的点。数量上两者只差有限个点，但共形块的定义
% 依赖的是挖掉点后的模空间。]

\SUB{4.8}{进一步推广}

在结束关于 Class S 构造的讨论之前，我们接着简要提及构造 4d $\cN=2$ SCFT 时手头
还有的几样尚未涉及的材料，即：6d 中 $\cN=(2,0)$ 理论的不同选择、Hitchin 场在每个
穿刺点附近的不同奇异行为，以及绕穿刺点一周时引入单值性扭曲（monodromy twist）的
可能性。

% [Tutor Note: 本段列出的三样「新原料」正好对应下面三个小节：换 6d $\cN=(2,0)$
% 理论的类型（A_{N-1} 之外还有 D_N、E_6、E_7、E_8）-> §4.8.1；允许 Hitchin 场有
% 更强或分数阶的极点（非正则穿刺点）-> §4.8.2；再加上绕穿刺点的外自同构单值性
% （扭曲线）-> §4.8.3。三者可以叠加使用，叠加后就得到 §5 几何工程要对照的那一大类
% 4d $\cN=2$ SCFT。]

\SUBSUB{4.8.1}{一般的 $\cN=(2,0)$ 理论}

在前面的各节中，我们主要关注使用一叠 $N$ 个 M5 膜的 Class S 构造，它描述
$A_{N-1}$ 型的 6d $\cN=(2,0)$ SCFT。还存在其他类型的 6d $\cN=(2,0)$ 理论，它们
分别是 $D_N$、$E_6$、$E_7$ 与 $E_8$ 型。$D_N$ 型理论可以在 M 理论中由一叠
$N$ 个 M5 膜放在一个 M 理论定向折叠面（orientifold）OM5 之上来构造 [223--225]。
更一般地，它们可以由 IIB 型弦理论探测 ADE 奇点来工程化，如 §4.1 所综述的那样。

% [Tutor Note: 为什么 D_N 型要点出 OM5？M 理论里没有弦耦合，不能像 IIB 那样用
% O7 平面；取而代之的是「M 理论定向折叠面」OM5。把一叠 M5 膜放在 OM5 上，等价于
% 用 Z_2 把 A 型链「折叠」，于是 A_{2N-1} 型被折成 D_N 型——这与 IIB 中 O7 平面
% 把 A 型折成 D 型是同一件事的 M 理论版本（见 [223-225]）。]

回忆在 $A_{N-1}$ 型 Class S 构造中，重要的 4d 库仑分支物理（例如手征主算子的谱、
手征耦合与味对称性）都包含在由 Higgs 场 $\Phi_z$ 构成的 $\fsu(N)$ 不变微分
$\phi_k$（$k=2,\dots,N$）之中。在抽象层面上，这些微分对应于 6d $\cN=(2,0)$ SCFT
中的半 BPS 算子 [5]；对一般的 ADE 型，它们与 Lie 代数 $\fg$ 的 Casimir 不变量
一一对应。对 $D_N$ 型 Class S 构造，独立的微分是 $k=2,4,\dots,2N-2$ 的 $\phi_k$
以及对应于 $\fso(2N)$ 的 Pfaffian 不变量的 $\tilde{\phi}_N$。对 $E_6$，微分是
$k=2,5,6,8,9,12$ 的 $\phi_k$。对 $E_7$，是 $k=2,6,8,10,12,14,18$；对 $E_8$，是
$k=2,8,12,14,18,20,24,30$。

% [Tutor Note: A_{N-1} 型的独立微分是 φ_k（k = 2, …, N），共 N-1 个，恰好等于
% su(N) 的秩；它们与 su(N) 的 Casimir 不变量（次数 2, 3, …, N）一一对应，也对应
% 4d 库仑分支上的 N-1 个库仑分支算子，且（正则穿刺点时）Δ[u_k] = k。原文随后列出
% 的 D_N、E_6、E_7、E_8 的次数表就是相应 ADE 代数的 Casimir 次数（例如 E_8 的
% 2,8,12,14,18,20,24,30 正是 E_8 的 8 个 Casimir 次数），可直接查表核对。]

% 说明：原文把表 4 作为浮动体排在 PDF p.70 顶端，恰在本段「…Higgs 场 Φ_z。At」与
% 下一句「the abstract level, …」之间（跨页处把一句话切开）。为不把一句话从中间
% 剖开，译文把本段补全后紧跟表 4，随后的「一大类…」段与原文顺序一致。

% [原文 p.69]
\begin{center}
\small
\begin{tabular}{>{\centering\arraybackslash}p{0.40\linewidth}
                >{\centering\arraybackslash}p{0.54\linewidth}}
\toprule
$\cN=2$ 理论 $\cT(A_1,\cC,p)$ & $\bar{\cC}$ 上的 Liouville 理论 \\
\midrule
形变参数 $\epsilon_1$ 与 $\epsilon_2$
  & Liouville 参数 $\epsilon_1=b$ 与 $\epsilon_2=b^{-1}$，以及
    $c=1+6Q^2$ 与 $Q=b+b^{-1}$ \\
四个自由超多重态 & 三穿孔的球面 \\
与某个 $\SU(2)$ 味相关联的质量参数 $m$ & 插入一个 Liouville 顶点算子 $e^{2m\phi}$ \\
带 UV 耦合 $\tau$ 的 $\SU(2)$ 规范群 & 粘合参数为 $q=\ee^{2\pi\ii\tau}$ 的 OPE 道 \\
某个 $\SU(2)$ 规范群的 VEV $a$
  & 中间道的主算子 $e^{2\alpha\phi}$，其中 $\alpha=Q/2+a$ \\
$Z_{\mathrm{inst}}(a,Q)$ & Virasoro 共形块 \\
$Z_{1\text{-loop}}(a)$ & OPE 系数的乘积（由 DOZZ 公式给出） \\
$Z_{S^4_b}(\cT(A_1,\cC,p))$ & Liouville 关联函数 \\
\bottomrule
\end{tabular}
\end{center}
\tabtitle{4}{秩一 AGT 对应中 2d 与 4d 可观测量之间的字典 [214]。}

% [Tutor Note: 表 4 逐行对应（左列 4d，右列 2d）：(1) 形变参数 ε_1, ε_2 <-> Liouville
% 参数 b, 1/b 与中心荷 c = 1 + 6Q^2（背景荷 Q = b + 1/b）；(2) 四个自由超多重态 <->
% 三穿孔球面；(3) 质量 m <-> 顶点算子 e^{2mφ}；(4) SU(2) 规范群与 UV 耦合 τ <->
% 粘合参数 q = e^{2πiτ} 的 OPE 道；(5) VEV a <-> 中间道主算子 e^{2αφ}，α = Q/2 + a；
% (6) Z_inst(a,Q) <-> Virasoro 共形块；(7) Z_{1-loop}(a) <-> OPE 系数（DOZZ 公式）；
% (8) 整球配分函数 <-> Liouville 关联函数。关于 (5)：Liouville 主算子的共形权为
% Δ = α(Q - α)，α 与 a 相差背景荷 Q/2，这正是 4d 库仑分支参数 a 与 2d 动量 α 的
% 平移关系。]
% [Tutor Note: 表 4 左列第 6 行原文印作 Z_inst(a, Q)，与 (4.27) 的 Z_inst(a, q) 不一致。
% 已用 800 dpi 裁图复核，原表确为大写 Q；按物理含义（瞬子计数参数是 q = e^{2πiτ}，
% 而 Q = b + 1/b 是 Liouville 背景荷）判断，此处应为 q。译文照排 Q，此处存疑。]

一大类 4d $\cN=2$ 理论可以由一个以 ADE Lie 代数标记的 6d $\cN=(2,0)$ SCFT
构造出来：把该 6d 理论作扭曲紧化，紧化在一条带有一族穿刺点的 Gaiotto 曲线 $\cC$ 上，
% [原文 p.70]
并聚焦到低能极限。要使所得的 4d 理论是 SCFT，$\cC$ 有如下几种选择，
\begin{itemize}[leftmargin=1.8em,itemsep=2pt]
  \item 带三个或更多正则穿刺点的球面。\footnote{由带三个正则穿刺点的球面构造出的
        SCFT 被称为 tinkertoy，并已在 [36,37,41,44--46,49,53,226] 中被广泛研究。}
  \item 带一个正则穿刺点和一个非正则穿刺点的球面。\footnote{非正则穿刺点将在
        §4.8.2 中引入。}
  \item 带一个非正则穿刺点的球面。
  \item 带正则穿刺点的高亏格 $\cC$。
\end{itemize}

% [Tutor Note: 这四个选项的物理含义：第一种（>=3 个正则穿刺点）就是 Gaiotto 的
% 原始构造；(2)(3) 需要非正则穿刺点，是 Argyres-Douglas 型理论的来源；第四种
% （高亏格 + 正则穿刺点）对应 N=4 类的理论（如单穿孔环面给 N=4 SU(2) SYM）。
% 注意原文第二项印作「one regular puncture and one irregular punctures」，复数
% punctures 与 one 不一致（已 500+ dpi 复核，原文如此），系行文笔误。]

\SUBSUB{4.8.2}{非正则（野）穿刺点}

原先在 [33] 中的 Class S 构造被极大地推广了，办法是纳入更一般的穿刺点，
即所谓的\textbf{非正则（野）}奇点。事实上，它们对于给出 Argyres--Douglas 型理论的
Class S 构造至关重要，而这类理论若只用正则穿刺点则是无法实现的。\footnote{带
未扭曲正则穿刺点的 Class S 构造有一个特征：所得 SCFT 的库仑分支手征环，其标度
维数谱全为整数。这一特征排除了大多数 Argyres--Douglas 型理论，因为后者的库仑分支
算子谱通常含有分数维数。不过，带扭曲正则穿刺点的 Class S 构造可以给出分数维数的
库仑分支算子 [227]。}

虽然对正则穿刺点而言 Hitchin 场只有单极点，但也存在一些穿刺点，其 Hitchin 场的
奇异性比式 (4.12) 中的更强。\footnote{它们定义 6d $\cN=(2,0)$ SCFT 中更一般的
余维二缺陷。}此外，Hitchin 场还可以有某些分数阶的极点，只要它们与结构群 $G$
相容 [35]。作为后果，$\phi_k$ 微分的极点阶数可以大于 $k$。我们称这类穿刺点为
\textbf{非正则（或野）}穿刺点，它们已在 [35,40,42,47,56,228,229] 中被研究过。
\footnote{请记住，我们这里使用的「非正则」穿刺点这一说法，指的是 Hitchin 场的
行为，而不是 [36] 中引入的定义。}
正则穿刺点有简单的分类：用标记 6d $\cN=(2,0)$ SCFT 的 Lie 代数 $\fg$ 的幂零轨道
来分类（例如经典 Lie 代数用杨图，更一般的情形用 Bala--Carter 标签 [41]）；相比之下，
非正则穿刺点的分类要丰富得多，并且存在一套系统的方法，可以从 Hitchin 极点读出
物理信息，例如味对称性（的 Cartan 生成元）、恰当边缘耦合以及库仑分支的手征主算子
（例如参见 [42,47]）。

% [Tutor Note: 正则穿刺点的分类是「离散且有限」的：幂零轨道（对 su(N) 就是杨图/
% 整数分拆）本身只依赖组合数据。而非正则穿刺点除了极点阶数（可含半整数）还要指定
% 极点展开的各阶系数，可调参数更多、分类「丰富得多」。正是这些额外参数让库仑分支
% 算子的维数可以取分数（见后面的 (A_1,A_2) 与 6/5）。]

% [Tutor Note: 「极点阶数可以大于 k」是理解非正则穿刺点为何能给出分数维算子的关键：
% 正则穿刺点时 φ_k 的极点阶数恒为 k，对应的库仑分支算子维数 Δ[φ_k] = k（整数），
% 这正是脚注 47 所说「未扭曲正则穿刺点的 CB 维数全为整数」。一旦允许分数阶极点，
% 由 SW 微分读出的维数就会出现分数（例 4.8 的 (A_1,A_2) 给出 Δ[u] = 6/5）。]
% [Tutor Note: 「与结构群 G 相容的分数阶极点」指的是：Φ_z 的极点阶数不能取任意实数，
% 只能取那些使留数落在 G 的相应子代数里的特定分数阶。在 A_1 情形允许的是半整数阶
% （如 z^{3/2}、z^{7/2}），这正是 (4.29) 与 (4.33) 中指数带 /2 的原因。]

